\documentclass{edpsci}
\usepackage{amsmath,amssymb,graphicx,booktabs,microtype}
\usepackage{etoolbox,balance}
\usepackage[numbers,sort&compress]{natbib}
\usepackage{hyperref}
\usepackage{url}
\urlstyle{rm}
\graphicspath{{figures/}}
\hypersetup{colorlinks=true,citecolor=blue,urlcolor=blue,linkcolor=blue}
\journalname{Acta Acustica}
\renewcommand{\submittedto}{Template preview for}
\articletype{Research article -- preprint}
\hypersetup{pdftitle={Effects of arrival timing, source startup and room reflections on three-microphone sound localization},pdfauthor={Aneesh Arnav Chikkala},pdfsubject={Acta Acustica template preview; unreviewed preprint; not submitted}}
% Keep the supplied class intact; accommodate its header and one unaffiliated author.
\makeatletter
\patchcmd{\@maketitle}{\vbox to 22.7mm}{\vbox to 27mm}{}{\errmessage{EDP title-box patch failed}}
\patchcmd{\makeheadbox}{\vbox to0pt}{\vbox}{}{\errmessage{EDP article-band patch failed}}
\renewcommand{\institutename}{\par\begingroup\parindent=0pt\noindent\@institute\par\endgroup}
\makeatother
% The class uses \doi for article metadata; bibliography entries need a printed DOI.
\renewcommand{\doi}[1]{doi: \href{https://doi.org/#1}{\nolinkurl{#1}}}

\begin{document}
\title{Effects of arrival timing, source startup and room reflections on three-microphone sound localization}
\titlerunning{Arrival timing and source startup in sound localization}
\author{Aneesh Arnav Chikkala}
\authorrunning{Aneesh Arnav Chikkala}
\institute{No institutional affiliation}
\date{8 September 2026}
\abstract{Real-time sound-source localization is essential for assistive hearing systems, wearable interfaces, robots, and distributed IoT sensors. However, the computational, memory, and energy requirements of conventional localization methods, such as beamforming, limit their use in low-cost, battery-powered edge devices. These constraints motivate a deliberately minimal architecture comprising three microphones and lightweight GCC-PHAT time-delay estimation. This study investigates whether temporal averaging can improve the direction-of-arrival (DoA) accuracy of this architecture without increasing per-frame processing, while examining how simulation and observation conditions affect the measured performance. Independent numerical calculations were used to verify the direct-path delays and simulated room responses. Matched simulations then examined the effects of arrival-time rounding, source onset, and reverberation. Separate recordings provided the data for assessing temporal averaging. At an SNR of 10~dB relative to the direct sound, rounding arrival times to the nearest sample increased the direct-path root-mean-square error (RMSE) from \(0.92^\circ\) to \(2.30^\circ\). In the more reverberant cases, the first frame after source onset gave an RMSE of \(4.16^\circ\), compared with \(42.51^\circ\) during steady operation. The lower error at source onset occurs because the reflected sound field has not yet fully developed. Averaging 128 frames, corresponding to 5.46~s, reduced the RMSE to \(0.84^\circ\) under lower reverberation and \(10.16^\circ\) under stronger reverberation. These results demonstrate that temporal averaging can substantially improve the accuracy of a computationally minimal DoA system, although the achievable performance remains limited by the acoustic environment. The findings also highlight the importance of controlling arrival-time discretization, source timing, reverberation, and observation duration when evaluating compact localization systems.}
\keywords{Sound localization, microphone arrays, GCC-PHAT, room simulation, arrival timing, source startup}
\maketitle
\raggedbottom

\section{Introduction}
Sound localization with a small microphone array is difficult because its limited aperture produces short time differences between microphone signals. These delays must be estimated in the presence of noise and reflections, which can obscure the direct sound. The present system uses three microphones and GCC-PHAT delay estimation, requiring the array geometry and processing method to be considered together.
At a sampling rate of 48 kHz, consecutive samples are $20.83,\mu$s apart, during which sound travels about 7.15 mm at 343 m/s. At this instance the microphone spacing is 86.6 mm, rounding an arrival time to the nearest sample introduces a significant error relative to the propagation time across the array. Arrivals between samples must therefore be represented accurately.
Simulation allows the effects of delay rounding, source onset, and reverberation to be examined separately. However, beginning the analysis when the source starts can produce misleading results. The direct sound reaches the microphones before the reflected field has fully developed, so the first frame may appear more accurate than measurements made during steady operation.
Generalized cross-correlation is an established method for estimating time delay \citep{knapp}, and GCC-PHAT has been applied to equilateral three-microphone arrays \citep{catur}. Fractional-delay filtering represents arrivals between samples \citep{laakso}, while the image-source method represents room reflections using mirrored sources \citep{allen}. The effect of these simulation and observation choices on the reported accuracy of small arrays has received less attention.
This study quantifies the effects of arrival-time rounding, source onset, and reverberation, and assesses temporal averaging using separate recordings. Direct-path delays and room responses are checked using independent numerical calculations, following the distinction between implementation verification and physical validation \citep{ernoult}. The simulations use one stationary synthetic source, three ideal matched microphones and three rectangular rooms, with localization restricted to azimuth. Microphone calibration, device timing and measured-room performance require separate experiments.

\section{Simulation design and implementation}

\subsection{Microphone array and direction estimation}

The three microphones form an equilateral triangle in the horizontal plane. Each microphone is 5 cm from the array centre, giving an inter-microphone distance of 8.66 cm. The source is located 1.5 m from the centre at the same height. Figure~\ref{fig:geometry} shows the geometry and the matched comparisons used in the study.

With azimuth measured counterclockwise from the positive $x$-axis, the microphone coordinates are
\begin{equation}
	\mathbf{m}_i
	=
	R[\cos\phi_i,\sin\phi_i]^T,
	\qquad
	\phi_i\in\{90^\circ,210^\circ,330^\circ\},
\end{equation}
where $R=0.05$ m. The speed of sound is fixed at $c=343$ m/s. For a source at $\mathbf{s}$, the direct-path arrival time at microphone $i$ is
\begin{equation}
	t_i=\frac{\lVert\mathbf{s}-\mathbf{m}_i\rVert}{c}.
\end{equation}
The time difference for pair $i,j$ is
\begin{equation}
	\tau_{ij}=t_i-t_j,\qquad i<j.
\end{equation}
The delay vector is ordered as
$\boldsymbol{\tau}=[\tau_{12},\tau_{13},\tau_{23}]^T$.
Although three pairs are available, their errors are correlated because the pairs share microphones.

Direction is estimated using a plane-wave model. For a unit vector $\mathbf{u}$ pointing toward the source,
\begin{equation}
	A\mathbf{u}=-c\boldsymbol{\tau},
\end{equation}
where each row of $A$ is $(\mathbf{m}_i-\mathbf{m}_j)^T$, ordered consistently with $\boldsymbol{\tau}$. The estimated delays give
\begin{equation}
	\widetilde{\mathbf{u}}
	=
	-cA^+\widehat{\boldsymbol{\tau}},
	\qquad
	\widehat{\theta}
	=
	\operatorname{atan2}
	(\widetilde{u}_y,\widetilde{u}_x),
	\label{eq:solve}
\end{equation}
where $A^+$ is the pseudoinverse. This is an equally weighted least-squares estimate; it is used as a fixed estimator rather than as a statistically optimal treatment of the correlated pairwise errors. The vector is not normalized because only its angle is used.

The finite source distance produces curved rather than planar wavefronts. Consequently, even exact direct-path delays can give a non-zero angular error when interpreted using equation~\ref{eq:solve}. This modelling error is quantified separately using the analytical direct-path arrival times.

\begin{figure*}[t]
	\centering
	\includegraphics[width=\textwidth]{figures/fig1_geometry_controls.pdf}
	\caption{Microphone geometry and matched comparisons. The triangle is drawn to scale, whereas the arrow indicates source direction only. The source is 1.5 m from the array centre. For each recording, matched conditions use the same source and microphone-noise samples. ``Onset'' denotes a single source startup at the beginning of the recording.}
	\label{fig:geometry}
\end{figure*}

\subsection{Room model and arrival timing}

Each room is represented as a three-dimensional rectangular enclosure with specular reflections. The six walls have frequency-independent pressure-reflection coefficients. Geometrical spreading is included, while diffraction, diffuse scattering, air attenuation and microphone directivity are excluded.

For room dimension $L_a$ and source coordinate $s_a$ along axis $a$, an image-source coordinate is
\begin{equation}
	s_a^{\mathrm{img}}
	=
	2n_aL_a+(1-2p_a)s_a,
\end{equation}
where $n_a\in[-K,K]$ and $p_a\in\{0,1\}$. The numbers of reflections at the low and high walls are $\lvert n_a-p_a\rvert$ and $\lvert n_a\rvert$, respectively. A path of length $r$ has an arrival time $r/c$ and pressure gain
\begin{equation}
	g(r)
	=
	\frac{1}{4\pi r}
	\prod_w\beta_w^{N_w},
\end{equation}
where $\beta_w$ and $N_w$ are the pressure-reflection coefficient and reflection count for wall $w$. The final model uses $K=76$ and retains paths up to 1.2 s. Convergence is assessed by repeating the room-response calculation with larger image extents and longer path limits.

For uniform walls, absorption is assigned using
\begin{equation}
	\alpha=\frac{0.161V}{S T_{\mathrm{nom}}},
	\qquad
	\beta=\sqrt{1-\alpha},
\end{equation}
where $V$ is the room volume, $S$ is its surface area and $T_{\mathrm{nom}}$ is the nominal reverberation setting. The resulting decay time is measured from the simulated response rather than assumed to equal $T_{\mathrm{nom}}$ \citep{bellows}.

Two arrival-time implementations are compared. In the fractional implementation, each path is deposited linearly on a grid 32 times finer than the acquisition grid. A normalized Kaiser-windowed sinc filter then converts the response to the acquisition rate. The Kaiser shape parameter is 8.6, and the filter half-width is 32 acquisition samples. In the rounded implementation, each path is placed at
\begin{equation}
	n=\operatorname{round}(f_s r/c).
\end{equation}
The path gains, reflection set and subsequent signal processing are unchanged. The comparison therefore measures the effect of replacing fractional arrival construction with direct sample rounding. A direct-only condition removes all reflections while retaining the same direct path and response length.

\subsection{Source signal and recording procedure}

The source is a Gaussian random sequence restricted to 300--3400 Hz using a Fourier-domain mask and normalized by its root mean square amplitude. The sampling rate is 48 kHz. Each continuous recording includes 59,682 source samples before the analysed interval, allowing the room response and causal acquisition filter to reach steady operation.

All microphone signals pass through the same causal Butterworth bandpass filter covering 300--3400 Hz. The SciPy prototype order is four, producing four second-order sections and a total bandpass order of eight. Independent Gaussian noise is generated for each microphone and passed through the same filter.

Noise is scaled relative to the filtered, steady-state, direct-only signal. Let $P_{\mathrm{dir}}$ denote the direct-signal power averaged over the microphones and analysed samples, and let $P_{\mathrm{noise}}$ denote the corresponding acquired noise power. The reported SNR is
\begin{equation}
	\mathrm{SNR}
	=
	10\log_{10}
	\left(
	\frac{P_{\mathrm{dir}}}{P_{\mathrm{noise}}}
	\right).
\end{equation}
The same noise samples and scaling are retained when reflections are added. The stated SNR is therefore referenced to the direct sound and is not the total reverberant-signal-to-noise ratio.

Source history is examined using matched steady and onset recordings. In the steady condition, the source is active during the pre-analysis interval. In the onset condition, the earlier source samples are set to zero once, while the source samples within the analysed interval and all microphone-noise samples remain unchanged. This comparison includes the effects of room-response buildup and causal-filter startup and is therefore treated as a comparison of recording conditions rather than as a pure reverberation comparison. Direct-only results are used to identify effects that remain when reflections are absent.

Each short recording contains 32 consecutive, non-overlapping frames of 2048 samples. A frame lasts 42.667 ms, and the complete short recording lasts 1.365 s. The temporal-averaging analysis uses a separate continuous recording containing 128 frames, corresponding to 5.46 s. For $N$ frame estimates, azimuth is averaged circularly:
\begin{equation}
	\overline{\theta}_N
	=
	\operatorname{atan2}
	\left(
	\sum_{n=1}^{N}\sin\widehat{\theta}_n,
	\sum_{n=1}^{N}\cos\widehat{\theta}_n
	\right).
\end{equation}
This avoids discontinuities at the azimuth wraparound boundary. A separate exploratory condition restarts the source and acquisition history before every frame. Because that procedure repeatedly samples the startup transient rather than continuous operation, its results are reported separately.

\subsection{Delay estimation using GCC-PHAT}

GCC-PHAT estimates the relative delay of each microphone pair from the peak of a phase-weighted cross-correlation. The mean is removed from each frame before applying a symmetric Hann window and a 4096-point Fourier transform. For microphone spectra $X_i[k]$ and $X_j[k]$, the cross spectrum is
\begin{equation}
	G_{ij}[k]=X_i[k]X_j^*[k].
\end{equation}
A denominator floor is applied before PHAT normalization:
\begin{align}
	d_{ij}[k]
	&=
	\max\bigl(\lvert G_{ij}[k]\rvert,q_{ij}\bigr),\\
	q_{ij}
	&=
	10^{-3}\max_l\lvert G_{ij}[l]\rvert+10^{-12},\\
	W_{ij}[k]
	&=
	G_{ij}[k]/d_{ij}[k].
\end{align}
The maximum in $q_{ij}$ is evaluated over all frequency bins. Values of $W_{ij}$ outside 300--3400 Hz are set to zero. The floor regularizes the PHAT weighting when the cross-spectrum magnitude is close to zero.

The inverse transform is evaluated on a correlation grid sampled eight times more finely than the acquisition grid. This interpolates the discrete correlation function but does not add information to the microphone signals. The peak search is restricted to
\begin{equation}
	\lvert\tau_{ij}\rvert
	\leq
	\frac{\lVert\mathbf{m}_i-\mathbf{m}_j\rVert}{c}.
\end{equation}
A quadratic fit to the peak and its two neighbouring samples provides a final sub-grid estimate. The correction is limited to one correlation-grid interval and clipped to the physical delay bound. The three pairwise estimates are then inserted into equation~\ref{eq:solve}.

\subsection{Measuring the simulated room response}

Room decay is calculated from the reverse-integrated squared impulse response using Schroeder's method \citep{schroeder}. Energy is summed across the three microphones before integration. A line is fitted to the decay between $-5$ and $-25$ dB and extrapolated to 60 dB to obtain T20. Early decay time, T30 and the goodness of fit are retained to assess whether the fitted range is supported by the truncated response.

The direct-to-reverberant ratio is calculated from separately filtered direct and reflected responses:
\begin{equation}
	\mathrm{DRR}
	=
	10\log_{10}
	\left(
	\frac{E_{\mathrm{dir}}}{E_{\mathrm{refl}}}
	\right).
\end{equation}
The interference term between the filtered responses is also retained because the energy of their sum is
\begin{equation}
	E_{\mathrm{tot}}
	=
	E_{\mathrm{dir}}
	+
	E_{\mathrm{refl}}
	+
	2\sum_n h_{\mathrm{dir}}[n]h_{\mathrm{refl}}[n].
\end{equation}
For each scene, the arrival times and gains of the direct path and six first-order reflections are recorded. These quantities are used to check the room construction and to interpret differences between nominal and measured decay.

\section{Quality control and evaluation}

\subsection{Independent numerical checks}

The direction calculation is first tested without room reflections. An independent scalar implementation solves the two-dimensional normal equations without calling the main direction-estimation routine. It tests plane-wave and spherical-wave arrivals, the cardinal directions, and all six permutations of the microphone channels. The analysis covers 720 source angles, array radii of 0.025, 0.05 and 0.10 m, source distances of 0.5, 1.5, 5 and 20 m, and sampling rates of 16, 48, 96 and 192 kHz. These combinations produce 34,560 geometric and sampling conditions. Both arrival models and all six channel permutations are evaluated within each condition.

A separate waveform calculation evaluates continuous-time harmonic signals and compares their expected phases with the output of the fractional-delay implementation. The signal bandwidth, recording duration and observation interval are kept identical. This separates errors in signal construction from errors in converting the estimated delays into an azimuth.

The room model is checked using an independent image-source enumeration in which every path is evaluated using its exact frequency-domain phase. This reference calculation does not use the main path generator or its interpolation filter. Path counts, spectra and estimated directions are compared for rooms with uniform and unequal wall coefficients. Image extent is increased until two successive calculations differ by less than $0.05^\circ$ in direction, 1\% in T20 and 0.1\% in the probe-spectrum comparison. Response duration and interpolation resolution are then increased separately and jointly to check that truncation and interpolation do not determine the result.

The $0.05^\circ$ criterion is a numerical convergence target used when interpreting effects of approximately $0.5^\circ$ or larger. It is not an accuracy requirement for the acoustic device. Because GCC-PHAT is a peak-based estimator, convergence of the room response does not guarantee that every random signal realization will produce the same selected peak. The pilot and main evaluations therefore use separate, reproducible random-number seeds.

\subsection{Test conditions and matched comparisons}

The three rooms measure $6\times5\times3$ m, $5\times4\times2.8$ m and $8\times6\times3.2$ m. The array centre is placed at coordinate fractions $(0.5,0.5,0.5)$ and $(0.6,0.45,0.5)$ of each room. Nominal reverberation settings of 0.15, 0.30 and 0.60 s are used. The 12 source directions are
\begin{equation}
	\begin{aligned}
		\theta_j &= 7.37^\circ+30^\circ j,\\
		j &= 0,\ldots,11.
	\end{aligned}
\end{equation}
The minimum distance between any source or microphone and a wall is 0.312 m.

The main evaluation contains all 216 combinations of room, array placement, nominal reverberation setting and source direction. Each scene is evaluated at SNRs of 0, 10 and 20 dB using five independent 32-frame simulated realizations. The resulting numbers of frame-level estimates are listed in Table~\ref{tab:protocol}. The allocation of conditions and the analysis rules were fixed before the held-out results were examined.

\begin{table}[t]
	\caption{Evaluation counts. Each frame contains 2048 samples. Pilot, reference and sensitivity calculations are additional.}
	\label{tab:protocol}
	\centering
	\small
	\begin{tabular}{@{}lrr@{}}
		\toprule
		Study & Base scenes & Estimates\\
		\midrule
		Main room study & 216 & 103,680\\
		Seven additional variants & 72 & 80,640\\
		Long realizations & 72 & 92,160\\
		\midrule
		Total & & 276,480\\
		\bottomrule
	\end{tabular}
\end{table}

Two subsets are used for matched comparisons. The low-reverberation/centred subset combines the centred array placement with the nominal 0.15-s setting. The high-reverberation/offset subset combines the offset placement with the nominal 0.60-s setting. Applying these conditions to all three rooms and 12 source directions gives 72 base scenes. Because both absorption and placement differ, comparisons between the subsets describe combined acoustic conditions and do not isolate reverberation time.

At 10 dB SNR, the same five source and noise realizations are used for all eight combinations of fractional or rounded arrivals, steady or onset history, and direct-only or reflected propagation. The fractional/steady/reflected condition is reused from the main evaluation only when all inputs are identical. For each scene, five additional 128-frame realizations are used to estimate bias and temporal dependence. Five separate 128-frame realizations form a held-out set for evaluating temporal averaging. These long realizations are independent of the pilot study and the 32-frame evaluation.

\subsection{Additional geometry and method checks}

A fixed subset contains four source directions in each of the six room conditions formed by the three rooms and the two matched acoustic conditions. These 24 scenes are used to examine a $17^\circ$ array rotation, an unequal triangular array, unequal wall absorption and a narrower 300--1700 Hz source band.

The acquisition filter and direction estimator retain their 300--3400 Hz passband in the narrower-source-band condition. Each variation uses five new 32-frame realizations at 10 dB SNR. Because these variations do not reuse the base-condition random inputs, they are treated as sensitivity analyses rather than matched estimates of individual causal effects.

Steered-response power with phase-transform weighting (SRP-PHAT) is included as a second localization method \citep{dibiase}. It uses the same signals, frame window, frequency band and PHAT normalization as GCC-PHAT. For pair $i,j$ and candidate azimuth $\theta$, define
\begin{equation}
	C_{ij}(\theta)
	=
	\sum_{k\in\mathcal{B}} W_{ij}[k]
	\exp\!\left[
	\mathrm{i}2\pi f_k\tau_{ij}(\theta)
	\right].
\end{equation}
The SRP-PHAT objective is then
\begin{equation}
	P(\theta)
	=
	\sum_{i<j}\operatorname{Re}
	\bigl\{C_{ij}(\theta)\bigr\}.
\end{equation}
Here $\mathcal{B}$ contains the frequency bins between 300 and 3400 Hz, and $\tau_{ij}(\theta)$ is the plane-wave delay predicted for pair $i,j$. The objective is evaluated on a $0.25^\circ$ azimuth grid and refined quadratically around its maximum. Results are checked against a $0.05^\circ$ grid. This comparison is restricted to the stated subset and is not intended as a general ranking of localization methods.

\subsection{Error measures and uncertainty}

For true azimuth $\theta$ and estimate $\widehat{\theta}$, angular error is defined as
\begin{equation}
	e=
	\operatorname{wrap}_{[-180^\circ,180^\circ)}
	\bigl(\widehat{\theta}-\theta\bigr).
\end{equation}
This wrapping prevents a small crossing of the angular boundary from being counted as a large error. Performance is summarized using root mean square error,
\begin{equation}
	\mathrm{RMSE}
	=
	\sqrt{
		\frac{1}{N}
		\sum_{n=1}^{N}e_n^2
	},
\end{equation}
together with median absolute error, the 90th percentile of absolute error, circular bias and the proportion of estimates with absolute error above $5^\circ$. The $5^\circ$ threshold was specified as a descriptive measure and is not a hearing or device-acceptance standard.

Uncertainty is estimated using 2000 bootstrap draws \citep{efron}. Within each scene, complete simulated realizations are sampled with replacement, and all frames from a selected realization remain together. The same bootstrap selections are applied to matched variants. This preserves the pairing between conditions and avoids treating frames from the same realization as independent experiments. The resulting intervals describe variation caused by the simulated source and microphone noise for the selected scenes. They do not describe uncertainty over all possible rooms, and the small number of independent realizations per scene limits the precision of the intervals.

Matched effects are calculated before aggregation. For a loss function $L$, the paired effect for realization $r$ is
\begin{equation}
	\Delta_r=L_{r,A}-L_{r,B}.
\end{equation}
Here $L$ is either signed angular error or squared angular error. Mean squared error is averaged before its square root is taken to obtain RMSE. A difference between two RMSE values is reported only as a difference in summary performance and is not interpreted as an independent error component.

Interactions are calculated as differences between paired effects. Let R and F denote rounded and fractional arrival construction, respectively. For propagation condition $q\in\{\mathrm{refl},\mathrm{dir}\}$, define the rounding effect as
\begin{equation}
	\Delta_{\mathrm{R}}^{(q)}
	=
	L_{\mathrm{R},q}-L_{\mathrm{F},q}.
\end{equation}
The rounding-by-reflection interaction is then
\begin{equation}
	\Delta_{\mathrm{int}}
	=
	\Delta_{\mathrm{R}}^{(\mathrm{refl})}
	-
	\Delta_{\mathrm{R}}^{(\mathrm{dir})}.
\end{equation}
This quantity measures whether the effect of arrival-time rounding changes when reflections are present. Per-scene results are retained alongside pooled summaries.

\subsection{Testing averaging on held-out realizations}

Temporal averaging is evaluated on the five held-out 128-frame realizations for each scene. The averaging lengths are
\begin{equation}
	N\in\{1,2,4,8,16,32,64,128\}.
\end{equation}
For each value of $N$, the first $N$ consecutive frame estimates from each realization are combined using the circular mean
\begin{equation}
	\overline{\theta}_N
	=
	\arg\!\left(
	\sum_{n=1}^{N}
	\exp\!\left[\mathrm{i}\widehat{\theta}_n\right]
	\right).
\end{equation}
The circular mean avoids discontinuities at the $-180^\circ/180^\circ$ boundary. The longest averaging interval contains 128 frames and lasts 5.461 s.

Let
\begin{equation}
	e_n
	=
	\operatorname{wrap}_{[-180^\circ,180^\circ)}
	\bigl(\widehat{\theta}_n-\theta\bigr)
\end{equation}
denote the frame error. When the errors are sufficiently concentrated, circular averaging can be approximated by averaging these scalar errors. Define the finite-lag covariance contribution as
\begin{equation}
	A_N
	=
	\sum_{\ell=1}^{N-1}(N-\ell)\gamma_\ell.
\end{equation}
Using this finite-lag covariance term, the variance and mean squared error are
\begin{equation}
	\begin{aligned}
		V_N
		&=
		\frac{\gamma_0}{N}
		+\frac{2A_N}{N^2},\\
		\operatorname{MSE}(\overline{e}_N)
		&\simeq b^2+V_N.
	\end{aligned}
	\label{eq:variance}
\end{equation}
Here $b$ is the persistent mean error, $\gamma_0$ is the frame-error variance and $\gamma_\ell$ is the covariance between errors separated by $\ell$ frames. The term $A_N$ collects the covariance contributions from non-zero lags. The bias term persists, whereas $V_N$ describes fluctuations that may decrease through averaging.

The bias and covariance terms are estimated separately for each scene using the five 128-frame estimation realizations. The persistent bias is obtained from their pooled circular mean error. Each realization is centred before its autocovariance is calculated, preventing between-realization offsets from being included as temporal covariance. Covariances above lag 32 are set to zero. For $1\leq\ell\leq32$, a Bartlett taper
\begin{equation}
	w_\ell=1-\frac{\ell}{33}
\end{equation}
is applied to reduce sampling variation at longer lags. A comparison model assumes independent frames by setting all non-zero-lag covariances to zero. Any small negative variance caused by finite-sample numerical error is set to zero before taking its square root.

Equation~\ref{eq:variance} is interpreted only for estimation-set results with a circular resultant of at least 0.95 and an absolute circular mean error no greater than $15^\circ$. These criteria were specified before examining the held-out realizations. They limit use of the small-error approximation but do not remove any scene from the measured held-out results. Scenes that fail the criteria remain in the empirical evaluation, and the criteria do not rule out rare large errors.

\section{Results}

\subsection{Numerical verification and correction of the room model}

Against the analytical solution, plane-wave recovery differed by at most $5.7\times10^{-14}$ degrees. The independent scalar solver differed from the main direction solver by at most $7.3\times10^{-12}$ degrees. The largest difference between the continuous-time waveform reference and the selected fractional-delay filter was $3.7\times10^{-5}$ degrees. Increasing the correlation-sampling factor from 8 to 64 changed the direct-path estimates by no more than $0.000495^\circ$.

The final independent room calculation evaluated about 5.22 million paths per microphone. Relative to this calculation, the main room model had a relative spectral difference of $1.13\times10^{-5}$ and an angular difference of $0.000807^\circ$. Earlier calculations with insufficient image coverage changed one direction estimate by more than $100^\circ$, and a coarse fractional-delay grid did not meet the spectral criterion. These failures were used to select the final numerical settings and are documented in the verification record.

A joint convergence check performed after the initial full run identified a further problem. Increasing response duration and image extent together changed ten of 960 checked estimates by more than $0.05^\circ$; the maximum change was $0.16551^\circ$. The original 0.9-s/extent-48 run was retained for audit but excluded from the results. The replacement model used a 1.2-s response, image extent 76 and propagation factor 32. Extending this model jointly to 1.5 s and extent 96 produced a maximum change of $0.01273^\circ$. The study was then rerun with the same allocation of conditions and the same analysis rules but with new source and noise identities. The original and corrected campaigns are therefore not treated as a realization-level matched comparison. All subsequent results are from the corrected run.

\begin{table}[t]
	\caption{Final numerical checks. Values are the largest changes in estimated angle under each check, rather than localization errors relative to the source.}
	\label{tab:verification}
	\centering
	\small
	\begin{tabular}{@{}lr@{}}
		\toprule
		Check & Change ($^\circ$)\\
		\midrule
		Independent full room reference & 0.000807\\
		Longer response and larger image extent & 0.012726\\
		Propagation grid 32 to 64 & 0.006965\\
		Image extent 76 to 84 & 0.000000\\
		GCC factor 8 to 16 & 0.025930\\
		SRP angular grid refinement & 0.002310\\
		\bottomrule
	\end{tabular}
\end{table}

Table~\ref{tab:verification} and Figure~\ref{fig:numerics} summarize the final checks. On the fixed 24-scene subset, doubling the propagation resolution changed the estimates by at most $0.006965^\circ$, increasing the image extent produced no change, and increasing the GCC correlation-sampling factor changed the estimates by at most $0.02593^\circ$. Refining the SRP angular grid changed the estimates by at most $0.002311^\circ$. Every checked change was below the $0.05^\circ$ numerical target. These results demonstrate convergence for the tested cases but do not provide a worst-case bound for every possible input.

\begin{figure*}[t]
	\centering
	\includegraphics[width=\textwidth]{fig2_numerical_verification.pdf}
	\caption{Arrival-time rounding and numerical resolution. (a) Angular change caused by rounding analytical arrivals for an array radius of 0.05 m and a source distance of 1.5 m. Exact spherical arrivals are processed by the same plane-wave direction solver, and the curves exclude waveform noise. (b) Largest angular changes under numerical refinement on the fixed 24-scene subset, using five recordings and 32 frames per scene. The dashed line marks the $0.05^\circ$ numerical target.}
	\label{fig:numerics}
\end{figure*}

The measured $T_{20}$ ranges were 0.083--0.110 s, 0.259--0.295 s and 0.586--0.683 s for nominal settings of 0.15, 0.30 and 0.60 s, respectively. The minimum fit $R^2$ was 0.928 for the lowest setting and 0.994 for the other settings. The component direct-to-reverberant ratio (DRR) ranged from $-10.77$ to $+9.53$ dB. Because the nominal settings specify the absorption configuration rather than the measured decay, both quantities are reported.

\subsection{Localization error across rooms}

At 10 dB SNR, the RMSE was $1.620^\circ$ for the nominal-0.15-s setting, $11.954^\circ$ for 0.30 s and $37.233^\circ$ for 0.60 s. The corresponding 95\% bootstrap intervals were $[1.599,1.641]^\circ$, $[11.168,12.771]^\circ$ and $[36.359,38.111]^\circ$. The proportions of estimates with absolute error above $5^\circ$ were 0.86\%, 27.38\% and 62.30\%, respectively.

At 20 dB SNR, the RMSE values were $1.362^\circ$, $11.796^\circ$ and $37.067^\circ$. The reduction in noise improved the lowest-reverberation result but had little effect on the larger errors in the more reflective conditions. Results for the full SNR range, including 0 dB, are shown separately for each room in Figure~\ref{fig:core}.

The observed scene-level RMSE at 10 dB ranged from $0.92^\circ$ to $3.16^\circ$ in the lowest setting, from $2.08^\circ$ to $49.08^\circ$ in the middle setting and from $6.56^\circ$ to $86.16^\circ$ in the highest setting. These ranges describe variation among the tested source directions and array placements. They should not be compared directly with the bootstrap intervals around pooled RMSE, which describe resampling uncertainty for the fixed collection of scenes.

\begin{figure*}[t]
	\centering
	\includegraphics[width=\textwidth]{fig3_core_rooms.pdf}
	\caption{Main room results obtained with fractional arrivals and a source that was already running. Each point combines two array placements, 12 source directions and five 32-frame recordings. Shaded bands are 95\% record-bootstrap intervals for the fixed scenes. The horizontal axis shows the nominal setting used to select the wall absorption; $T_{20}$ was measured separately.}
	\label{fig:core}
\end{figure*}

\subsection{Effect of rounding arrival times}

Across the 72 matched scenes at 10 dB SNR, the steady direct-path RMSE was $0.923^\circ$ $[0.912,0.934]^\circ$ with fractional arrivals and $2.300^\circ$ $[2.287,2.313]^\circ$ with rounded arrivals. The matched increase in mean squared error was $4.438\,\mathrm{deg}^2$ $[4.380,4.493]\,\mathrm{deg}^2$. The difference therefore remained after waveform construction, filtering and delay estimation; it was not confined to calculations based on exact delays.

With reflections included, the steady RMSE was $27.365^\circ$ for fractional arrivals and $28.112^\circ$ for rounded arrivals. The matched increase in mean squared error was $41.448\,\mathrm{deg}^2$ $[11.974,70.571]\,\mathrm{deg}^2$. This pooled result does not imply a uniform effect across scenes. Altering reflected-path arrival times can change which cross-correlation peak is selected, so both the magnitude and direction of the scene-level change can vary.

\subsection{Effect of source startup}

The strongest startup effect occurred in the first analyzed frame of the high-reverberation/offset subset. The RMSE was $4.157^\circ$ when the source began immediately before acquisition and $42.508^\circ$ when it was already running. By the final frame of the 1.365-s recording, both matched conditions had an RMSE of $41.526^\circ$. At that point, the earlier source history no longer contributed through the finite simulated response. Figure~\ref{fig:paired} shows the convergence of the two conditions as the reflected field develops.

The algorithm was unchanged between the onset and steady conditions; the difference arose from the sound field presented to it. Restarting the source before every short observation repeatedly samples the initial period before the reflected field has developed and therefore does not represent continuous operation. The exploratory frame-reset check produced the same qualitative behavior, but it used the pilot directions and is not included in the main time trajectory.

For fractional reflected signals across all 72 scenes and 32 frames, startup reduced mean squared error by $42.798\,\mathrm{deg}^2$, with a reduction interval of $[22.983,62.404]\,\mathrm{deg}^2$. Table~\ref{tab:effects} reports the same comparison as the signed onset-minus-steady contrast, for which negative values indicate less error at onset.

The interaction contrasts were calculated from the complete factorial dataset, not from the steady-only summaries above. The onset-by-reflection interaction was $-49.046\,\mathrm{deg}^2$ $[-65.574,-32.891]\,\mathrm{deg}^2$. The rounding-by-reflection interaction was $30.762\,\mathrm{deg}^2$ $[4.770,57.791]\,\mathrm{deg}^2$. Under the tested conditions, these non-zero interactions show that the effects of arrival construction, source history and reflected propagation are not additive.

\begin{table}[t]
	\caption{Matched changes in mean squared error ($\mathrm{deg}^2$) across 72 scenes at 10 dB. Positive values indicate more error. Intervals were obtained by resampling complete recordings within each scene.}
	\label{tab:effects}
	\centering
	\small
	\begin{tabular}{@{}lrr@{}}
		\toprule
		Comparison & Change & 95\% interval\\
		\midrule
		Rounding, direct & 4.44 & [4.38, 4.49]\\
		Rounding, reflected & 41.45 & [11.97, 70.57]\\
		Onset, reflected & -42.80 & [-62.40, -22.98]\\
		Reflections, fractional & 747.97 & [699.52, 799.81]\\
		\bottomrule
	\end{tabular}
\end{table}

\begin{figure*}[t]
	\centering
	\includegraphics[width=\textwidth]{fig4_paired_attribution.pdf}
	\caption{Arrival construction and source startup at 10 dB. (a) All eight combinations. F and Q denote fractional and rounded arrivals, S and O denote steady and onset histories, and D and R denote direct and reflected propagation. The low and high subsets contain 36 scenes each and differ in both array placement and wall absorption. Error bars were obtained by resampling complete recordings. (b) Matched fractional-arrival recordings with reflections in the high subset. In the steady condition, the source was already running when acquisition began. Markers indicate frame midpoints after the start of analysis.}
	\label{fig:paired}
\end{figure*}

\subsection{Improvement from temporal averaging}

Temporal averaging reduced block-level RMSE on the independent test recordings (Figure~\ref{fig:averaging}). In the low-reverberation/centred subset, RMSE fell from $1.597^\circ$ $[1.584,1.610]^\circ$ for one frame to $0.835^\circ$ $[0.815,0.854]^\circ$ for 128 frames. In the high-reverberation/offset subset, it fell from $39.044^\circ$ $[38.439,39.607]^\circ$ to $10.158^\circ$ $[9.404,10.995]^\circ$. These values come from the held-out averaging recordings and therefore need not equal the frame-specific values in the startup comparison. After 128 frames, 46.11\% of the high-subset block estimates still had an absolute error above $5^\circ$.

Forty-nine of the 72 scenes met the estimation-set criteria for interpreting the small-error approximation. Their measured 128-frame test RMSE was $2.013^\circ$ $[1.932,2.090]^\circ$. The prediction that included frame covariance was $2.075^\circ$, whereas the independent-frame prediction was $2.106^\circ$. These comparisons are descriptive because uncertainty intervals were not calculated for the model predictions. The full duration-dependent results are shown in Figure~\ref{fig:averaging}.

All 72 scenes remain in the empirical averaging curves. The restricted 49-scene analysis evaluates the approximation only on scenes that met its estimation-set criteria and should not be interpreted as performance for the complete test set. The errors remaining after 5.46 s varied among scenes and estimators; the finite observation period does not determine a universal or infinite-time reverberation limit.

\begin{figure*}[t]
	\centering
	\includegraphics[width=\textwidth]{fig5_heldout_averaging.pdf}
	\caption{Temporal averaging evaluated on new recordings. (a) Circular block averages for all 72 scenes, separated into the low-reverberation/centred and high-reverberation/offset subsets. (b) Results for the 49 scenes that met the criteria applied to the estimation recordings. Model curves are predictions, whereas shaded bands are bootstrap intervals obtained from complete test recordings. Estimation and test recordings are independent.}
	\label{fig:averaging}
\end{figure*}

\subsection{Timing-precision bound for the array}

An additional calculation relates errors in microphone arrival time to errors in the fitted direction. Let $B$ form the pairwise differences between the three microphone arrivals, and let $M$ contain the microphone coordinates, with $A=BM$. If the microphone arrival errors are collected in $\boldsymbol\epsilon$, the corresponding change in the fitted direction vector is
\begin{equation}
	\delta\widetilde{\mathbf u}
	=
	J\boldsymbol\epsilon,
	\qquad
	J=-cA^+B,
\end{equation}
where $A^+$ is the Moore--Penrose pseudoinverse. This expression retains the dependence introduced when the same microphone appears in more than one pair.

Suppose that each arrival error satisfies $|\epsilon_i|\leq h$. Because the norm of a linear mapping is convex over this error box, its maximum can be found from the eight corner combinations $\epsilon_i=\pm h$:
\begin{equation}
	\rho
	=
	\max_{\epsilon_i\in\{-h,h\}}
	\left\|J\boldsymbol\epsilon\right\|.
\end{equation}
If $\rho<\|\widetilde{\mathbf u}\|$, every admissible perturbation satisfies the conservative angular bound
\begin{equation}
	|\delta\theta|
	\leq
	\arcsin\!\left(
	\frac{\rho}{\|\widetilde{\mathbf u}\|}
	\right).
	\label{eq:bound}
\end{equation}
For the centred equilateral array and a plane wave, $\rho=4ch/(3R)$ and $\|\widetilde{\mathbf u}\|=1$. A sufficient condition for keeping the angular change below $\eta$ is therefore
\begin{equation}
	h
	\leq
	\frac{3R}{4c}\sin\eta.
\end{equation}
At $R=0.05$ m and $\eta=0.5^\circ$, this sufficient condition permits an absolute error of at most $0.954\,\mu$s in each arrival time. The result applies to ideal arrival times interpreted by the direction solver; it is neither a necessary sampling-rate requirement nor a guarantee for noisy GCC-PHAT estimates.

The numerical verification compared each observed angular change with its calculated bound. None of the 34,560 rounded-arrival cases exceeded the bound. A further 12,000 random perturbations, including unequal triangles and an independent scalar solution, also produced no violations. Figure~\ref{fig:precision}(a) measures each change relative to the estimate obtained from exact spherical arrivals. The difference between that reference estimate and the true source direction is a separate finite-distance effect.

\subsection{Additional geometries and comparison with SRP-PHAT}

On the fixed 24-scene subset, the GCC-PHAT RMSE was $22.413^\circ$ in the base condition. The independently simulated sensitivity conditions produced RMSE values of $23.729^\circ$ for the rotated array, $19.644^\circ$ for the unequal triangle, $23.643^\circ$ for unequal walls and $42.054^\circ$ for the narrower-band source (Figure~\ref{fig:precision}(b)). Because these conditions used new random recordings, their differences from the base value are not paired estimates of causal effects and do not identify an optimized geometry. In the narrower-band condition, part of the fixed estimator passband contained noise but no source energy; its result therefore reflects both the narrower source spectrum and this estimator-band mismatch.

The SRP-PHAT comparison used the same base recordings as GCC-PHAT. SRP-PHAT produced an RMSE of $30.262^\circ$, compared with $22.413^\circ$ for GCC-PHAT, but its proportion of errors above $5^\circ$ was lower: 28.62\% compared with 30.57\%. The matched SRP-minus-GCC difference in mean squared error was $413.473\,\mathrm{deg}^2$ $[268.788,561.435]\,\mathrm{deg}^2$. The two summaries describe different parts of the error distribution: fewer threshold exceedances can occur alongside a larger contribution from extreme errors. This restricted comparison does not establish a general ranking of the two methods.

\begin{figure*}[t]
	\centering
	\includegraphics[width=\textwidth]{fig6_precision_sensitivity.pdf}
	\caption{Timing precision and sensitivity conditions. (a) Analytical changes caused by arrival-time rounding and their bounds at 48 kHz, with an array radius of 0.05 m and a source distance of 1.5 m. Changes are measured relative to the estimates obtained from exact spherical arrivals. (b) RMSE and record-bootstrap intervals for five conditions on the fixed 24-scene subset. In the narrower-source-band condition, the acquisition and estimator bands remain 300--3400 Hz.}
	\label{fig:precision}
\end{figure*}

\section{Discussion and practical implications}

Subsample arrival construction materially affected the direction estimates. Fractional-delay methods are intended to represent arrivals that lie between discrete sample times \citep{laakso}; replacing those arrivals with nearest-sample values increased error even after waveform generation and GCC-PHAT processing \citep{knapp}. The effect was small relative to the largest reverberation errors but large enough to affect evaluations of compact arrays.

Source history had a larger influence on the first reflected-sound frame. In an image-source simulation \citep{allen}, setting the signal to zero before acquisition removes contributions that would otherwise arrive from earlier emissions. The resulting onset frame therefore represents a different acoustic condition from steady operation. A continuous-operation simulation requires sufficient prehistory to cover both the retained room response and the acquisition-filter transient. Conversely, an application concerned specifically with source onset should retain the onset condition and label it accordingly.

The numerical correction shows why response duration and image extent must be checked together. A long response window cannot include paths that were never generated because the image extent was too small. Agreement between response spectra is also insufficient when a peak-based estimator can select a different correlation maximum. The joint response and direction checks identified the limitation in the original room calculation and motivated the corrected campaign.

Temporal averaging reduced the measured error without increasing the processing applied to each frame. Its effect depended strongly on the scene. Averaging can suppress frame-to-frame fluctuations, but it cannot remove wavefront-model error, array-geometry limitations or persistent selection of an incorrect correlation peak. The covariance model was assessed only on scenes that met the small-error criteria, while the empirical curves retained the complete set. This distinction prevents the approximation from being presented as a model of the difficult scenes that it excludes.

The matched SRP-PHAT comparison \citep{dibiase} also illustrates the importance of reporting more than RMSE. SRP-PHAT produced fewer errors above the selected threshold but a larger RMSE, indicating greater influence from large errors. Because the comparison used one restricted subset and fixed implementations of both estimators, it does not support a general claim that either method is superior.

Several limitations restrict physical interpretation. The microphones are ideal and matched, so the simulations omit gain mismatch, clock drift, calibration error and device-specific frequency response. The source is stationary band-limited Gaussian noise rather than speech or machinery. The room model uses frequency-independent, specular wall reflections and does not include measured scattering or directional microphone responses. The main study also uses one array size and one source distance, although the analytical timing checks cover a wider range. Only horizontal direction is estimated, and neither the source nor the array moves during a recording.

The bootstrap intervals are based on five independent short recordings per scene. They describe resampling variation for the fixed simulated scenes and should not be interpreted as calibrated population intervals for arbitrary rooms. Their coverage was not assessed independently, and scene-level intervals have limited resampling support. The three room geometries likewise do not represent the range of real rooms. A measured-room experiment is therefore required to assess transfer to physical microphones and real acoustic environments.

\section{Conclusion}

The simulations separate the effects of arrival-time rounding, source startup and reflected propagation on three-microphone direction estimation. Nearest-sample rounding increased direct-path error after waveform processing, while beginning acquisition at source onset produced a temporarily easier reflected-sound condition than steady operation. Temporal averaging reduced error on independent recordings, but the remaining error depended on the scene.

Independent calculations and refinement tests establish the numerical resolution of the reported comparisons for the tested cases. The geometry-based timing bound provides a separate check on ideal arrival perturbations while retaining the dependence between microphone pairs. The results define which simulation and observation conditions must be reported when evaluating a compact localization system. They do not establish hardware accuracy, universal estimator superiority or a general relation between reverberation and localization error.


\bibliographystyle{unsrtnat}
\balance
\begingroup\interlinepenalty=10000
\bibliography{references}
\endgroup
\end{document}